% 補足資料 S1 の表：組と乱数のシード，評価に用いたチェックポイント（第7弾，2026-10-28。第8弾 2026-10-29 に組4〜6 の符号無作為条件を学習済みに直した：data/result/paper_extra/C_signrand/s1_signrand_{mse,l1_s3..5} の ep200，作業記録 §34.8）。
% 出典：data/result/ の学習の置き場所（シードごとの番号と train_seed<s>.log），学習の起動のスクリプト（--seed），
%   config/exp の L1 期間・段階IIの設定（pretrain の seed・epoch），評価のスクリプト（run_eval_s1・run_s2eval・run_c_eval・b4_run・eval_inverse・eval_u の既定の
%   エポック 160 170 180 190 200），docs/paper_extra_log.md §2.2・§4.3 の4（チェックポイントの一覧）。照合の記録：paper/check/numbers_check_v7.md §2。
\begin{tabular}{@{}p{0.05\linewidth}p{0.06\linewidth}p{0.07\linewidth}p{0.42\linewidth}p{0.27\linewidth}@{}}
\toprule
組 & シード & 区別 & 段階Iで学習した条件 & 段階IIの構成 \\
\midrule
組1 & 0 & 本番 & 経験なし，経験50\%，経験80\%，符号無作為，零入力，非共有（経験80\%・経験なし），好み切替 & 自由出力，形制限，形制限の参考，他者遮蔽学習，1ビット \\
組2 & 1 & 本番 & 組1 と同じ & 自由出力，形制限，他者遮蔽学習，1ビット \\
組3 & 2 & 本番 & 組1 と同じ & 自由出力，形制限，他者遮蔽学習，1ビット \\
組4 & 3 & 予備 & 経験なし，経験50\%，経験80\%，符号無作為 & 自由出力，形制限 \\
組5 & 4 & 追加 & 経験なし，経験50\%，経験80\%，符号無作為 & ― \\
組6 & 5 & 追加 & 経験なし，経験50\%，経験80\%，符号無作為 & ― \\
\midrule
\multicolumn{5}{@{}l}{評価に用いたチェックポイント（エポックの番号．最終5点は 160，170，180，190，200）} \\
\midrule
\multicolumn{3}{@{}p{0.2\linewidth}}{自分側の指標} & \multicolumn{2}{p{0.7\linewidth}@{}}{段階Iの MSE 期間の最終5点（自己位置の追跡，価値と行動・色の結びつき，$g$，$d_P$，次の視界の予測誤差）} \\
\multicolumn{3}{@{}p{0.2\linewidth}}{他者側の指標} & \multicolumn{2}{p{0.7\linewidth}@{}}{段階Iの L1 期間の最終5点（G5a，G5b，符号に従う分）} \\
\multicolumn{3}{@{}p{0.2\linewidth}}{真値条件の誤差と $\Delta$} & \multicolumn{2}{p{0.7\linewidth}@{}}{段階Iの L1 期間の 200 の1点（他者側の特徴予測の誤差は L1 期間のどの点でも MSE 期間の 200 と同じ）} \\
\multicolumn{3}{@{}p{0.2\linewidth}}{$h^2$ からの A-2 の位置} & \multicolumn{2}{p{0.7\linewidth}@{}}{段階Iの MSE 期間の 200} \\
\multicolumn{3}{@{}p{0.2\linewidth}}{照合} & \multicolumn{2}{p{0.7\linewidth}@{}}{自己視界基準は L1 期間の最終5点，他者特徴基準は MSE 期間の最終5点} \\
\multicolumn{3}{@{}p{0.2\linewidth}}{段階IIの指標} & \multicolumn{2}{p{0.7\linewidth}@{}}{段階IIの最終5点（段階Iの部分は L1 期間の 200）} \\
\multicolumn{3}{@{}p{0.2\linewidth}}{再現の確認} & \multicolumn{2}{p{0.7\linewidth}@{}}{200 の点を2回流し，ビット単位の一致を確かめた} \\
\bottomrule
\multicolumn{5}{@{}p{0.98\linewidth}@{}}{注：シードは学習の起動時に与え（L1 期間は同じシードの MSE 期間の 200，段階IIは同じシードの L1 期間の 200 から始める），データは組によらず同じである．経験なし条件のモデルは，データの一致と，学習し直したときの重みのビット単位の一致を確かめたうえで，経験50\%条件と経験80\%条件の両方の相手として用いた．段階IIの自由出力・形制限・他者遮蔽学習は経験なし・経験50\%条件の上で（自由出力は経験80\%条件の上でも），緑好みデータと好み混合データで学習し，形制限の参考の構成は組1 だけで行った．1ビットモジュールは経験なし・経験50\%・経験80\%・好み切替条件の上で，好み混合データで学習した．符号無作為条件は組1〜6，零入力条件・非共有条件・好み切替条件は組1〜3 で学習した．} \\
\end{tabular}
